% Page 051 — page imprimée 48
% Contenu seul : le préambule, l'en-tête et la pagination sont gérés par meyrueis.tex

\begin{center}
{\Large\bfseries Les Chantiers de Jeunesse : le groupement 19 à Meyrueis\par}
\end{center}

Le premier hiver 1940-1941 fut très froid. Les 130 ou 140 jeunes par groupes
venaient de s’installer et certains le passèrent sous la tente, en particulier à Valbelle,
au dessus des Ouberts à 1.100m, à l’hubac, donc sans soleil. On peut espérer que les
moins aguerris furent hébergés de temps à autre à la Maison Forestière. C’était le
groupement des « disciplinaires » où étaient expédiées les fortes têtes.

\textbf{La Pépinière}

\emph{« Mais il y eut des groupes mieux lotis. Je raconterai très brièvement l’histoire de
celui de Rousses. Il avait été affecté à ce hameau déjà en mauvais état : il y avait
bien là quelques maisons inocupées mais bien à l’étroit au bord du Béthuzon pour y
installer tout le groupe ; or le chef de groupe avait découvert, un peu plus haut, un
endroit beaucoup plus favorable. L’idée était judicieuse, elle fut acceptée par le
commandement. Mais il fallait envisager la construction d’un pont. Quelques fortes
têtes faisaient partie de ce groupe ; leur chef leur dit : « j’ai besoin de durs, de gens
qui n’aient peur de rien. Il s’agit de faire un pont et d’aménager une route pour
gagner l’autre rive » Ils réagirent avec enthousiasme Et on mit au travail :
remise en état du chemin muletier sur la rive gauche du Béthuzon à partir du petit
pont de Rousses puis construction du pont juste avant le camp, où fut fixée une
pancarte le jour de l’inauguration « Le pont des durs ». Ainsi fut installé le groupe
de « la Pépinière ». »}

Non contents de ce travail, ces jeunes construisirent pour les chefs une salle à
manger à 2m. du sol dans un énorme fayard. On y grimpait par une échelle !! et sur
la lancée, sensibilisés par leur chef de groupe, ils voulurent installer eux-mêmes leur
électricité dans les baraques qu’ils avaient construites, en utilisant l’eau du
Béthuzon. Ils avaient découverts que, au Villaret me semble-t-il, les gens avaient
fabriqué leur électricité. Ils céderaient volontiers un peu de matériel. Un des chefs
du Chantier se débrouilla pour trouver une petite turbine : une roue « pelton » de 50
ou 60 cm de diamètre, bien suffisante. La turbine tournait à toute allure et les
baraques étaient illuminées. Quelle merveille !!! \emph{(J’ai pu voir encore le bassin de
charge et les restes du canal d’amenée d’eau en Octobre 2006)}

\textbf{La vie du chantier}

Le ravitaillement était un problème important et un souci de tous les jours en raison
des restrictions et de la dispersion des groupes.
Les transports étaient assurés par des camions à gazogène car l’essence, très rare ;
était destinée aux médecins, y compris les médecins du pays, aux aumôniers et aux
chefs de groupement. « Heureusement nous disposions de quelques bicyclettes et
chevaux mais surtout nous pouvions compter sur nos 120 mulets ».

Les journées commençaient avec le rassemblement, le lever aux couleurs et
l’éducation physique ; ensuite alternaient le travail forestier, des aménagements de la
vie matérielle et des baraques ou des foyers, on chantait beaucoup, semble-t-il, on
faisait du sport régulièrement ; feux de camp le soir.
Quelle vie idyllique d’après les souvenirs de quelques-uns !!!
